\section{The Economic and Structural Context of the Collapse}
\label{sec:context}

The statistical results locate the collapse in time and shape but not in cause; the
economic and structural history of Albania between the 2018 and 2022 assessments
supplies the missing narrative. PISA~2018 was fielded in early 2018 and PISA~2022 in
2022, so any explanation must point to forces that acted in that four-year window. Two
shocks did, and only one of them was shared with Albania's peers.

\paragraph{Two compounding shocks hit the schooling window.}
Albania entered the pandemic already weakened by a natural disaster that its comparison
countries did not face. On 26 November 2019 a magnitude~6.4 earthquake struck the
Durr\"es and Tirana region, the strongest to hit the country in more than four decades;
it killed 51 people, injured about 3{,}000, and damaged or destroyed housing and
schools across roughly a dozen cities in the most populous part of
Albania~\cite{albania_pdna2020}. Barely three months later, COVID-19 closed schools
nationwide from March~2020, and Albania recorded one of the longer cumulative closure
durations in the region before instruction fully
resumed~\cite{unesco2022closures,worldbank_covid_learning}. The comparison economies in
this study lived through the pandemic, but none of them absorbed a major earthquake in
the same window. Notably, the OECD's own country note attributes Albania's 2022 decline
to exactly this pairing, an end-of-2019 earthquake that damaged schools and houses
across twelve cities, compounded by the pandemic~\cite{oecd_albania_note2023}. Albania
therefore experienced a \emph{double} disruption to the continuity of schooling
precisely between the two assessments, which is where a structural break would have to
originate.

\paragraph{The shock was to schooling, not to household wealth, and the data show it.}
If the 2022 collapse were a household-income shock, the socioeconomic markers in the
questionnaire would have fallen; instead, the school-side marker fell while the
household marker held. The covariate-shift analysis (Section~\ref{sec:results}) found
that the variable which moved most from 2018 to 2022 was teacher support in
mathematics, down by roughly a third of a standard deviation, while reported home
possessions actually rose. A demand-side poverty shock would depress home possessions;
the observed pattern instead points to a supply-side degradation of instruction, the
signature of closed and damaged schools and of teaching capacity stretched thin. This
interpretation is consistent with the macroeconomic record: Albania's economy
contracted in 2020 and rebounded strongly the next year, with World Bank figures
placing the contraction at roughly three to four percent and 2021 growth near nine
percent as tourism recovered~\cite{worldbank_albania}, while household consumption was
cushioned throughout by diaspora remittances worth about nine percent of
GDP~\cite{worldbank_remittances}. Family assets did not crater, which is why home
possessions did not; the damage concentrated in the public delivery of schooling.

\paragraph{Structural weaknesses turned a temporary shock into a lasting shift.}
A closure need not become a permanent loss, but three structural features of Albania's
system made it likely to. First, remote learning depended on infrastructure that many
households lacked: home internet and device access, especially in rural areas, left a
substantial share of students unable to follow online instruction, so televised and
online lessons reached pupils unevenly~\cite{worldbank_covid_learning,unicef_albania}.
Second, Albania spends a comparatively low share of GDP on education, on the order of
three to four percent, below the world average of about 4.6 percent and below European
Union norms~\cite{worldbank_edu_spending}, leaving little slack to mount an intensive
catch-up response. Third, Albania has for years experienced heavy emigration of
working-age and skilled adults, including teachers; the 2023 census confirmed that the
resident population had fallen by roughly 14 percent, about 420{,}000 people, since
2011, to some 2.4 million~\cite{instat_census2023}. This brain drain thins instructional
capacity at precisely the moment recovery demands more of it. Together these conditions
explain why the disruption still registered, undiminished, in the 2022 scores rather
than bouncing back.

\paragraph{The economic reading matches the statistical shape of the crisis.}
A system-wide capacity shock predicts a broad-based decline, and that is what the
gradient analysis found. Because closures, damaged buildings, and thinned teaching
capacity act on all students in a school rather than only its poorest, such a shock
should depress even advantaged students and flatten the socioeconomic gradient. Albania
in 2022 has the flattest SES gradient of the nine countries and a floor so low that
even its advantaged students score near the proficiency line
(Section~\ref{sec:results}); its mean mathematics score fell from 437 in 2018 to 368 in
2022, a drop of about 69 points~\cite{oecd_albania_note2023}. A shock concentrated in
household poverty would have steepened the gradient, not flattened it. A system-wide
shock also predicts a \emph{geographically} broad decline, and a difference-in-differences
confirms it: the mathematics collapse is no larger in the earthquake-affected central
region than elsewhere once a placebo exposes diverging pre-trends
(Section~\ref{sec:results}), so the 2019 earthquake compounded but did not localise the
loss, the decisive disruption was the national, pandemic-era interruption of schooling.
The convergence of the economic history and the statistical shape is the central insight
of this section: the collapse behaves like a supply-side disruption to schooling, not a
demand-side collapse in family means.

\paragraph{A caveat on measured effort.}
One qualification tempers the magnitude, though not the direction, of the drop. The OECD
country note reports that Albanian students spent less effort on the 2022 assessment
than in 2018 and shows signs of lower test engagement~\cite{oecd_albania_note2023}, so
part of the measured score fall may reflect test-taking behaviour rather than learning
alone. This does not undo the structural reading. Our central shift diagnostic
(Section~\ref{sec:results}) is trained on background questionnaire responses, not on the
scores themselves, and still separates the 2018 and 2022 cohorts at AUC~0.98; the change
in \emph{who} the students are is real independent of how hard they tried on test day.

\paragraph{A caveat on sample coverage.}
A second qualification concerns who the 2022 sample represents. The OECD flags that one
or more of its sampling standards were not met for Albania in PISA~2022, and Albania's
Coverage Index~3, the share of the national 15-year-old population that the assessed
sample represents, is about 79\%: the 6{,}129 students tested in 274 schools stand in
for roughly 28{,}400 of an estimated 36{,}000
15-year-olds~\cite{oecd2024pisa,oecd_albania_note2023}. The out-of-frame fifth is
disproportionately made up of adolescents who have left school or were never enrolled,
who as a group would be expected to perform below those still in the system. If
anything, then, this exclusion biases the measured at-risk rate \emph{downward}, so the
true breadth of low proficiency is likely wider than the already severe 75\% we report,
and the collapse is not an artefact of sweeping in weaker students. We make this precise
with partial-identification bounds (Section~\ref{sec:results}): even the worst-case
Manski interval that makes no assumption about the uncovered fifth is
$[0.58,\,0.79]$, and the monotone bound that treats them as no better off than the
assessed is $[0.74,\,0.79]$, so under any defensible assumption the crisis rate stays
far above Albania's 2018 level of 42\%. Two design points
guard the central comparison further. Coverage was of comparable order in 2018, so the
2018-to-2022 contrast is not driven by a coverage change; and our structural-shift
diagnostic conditions on background questionnaire responses within the assessed
population rather than on population totals, so it is insensitive to the level of
coverage even if it shifted. Coverage does, however, counsel caution in reading the
absolute rate as a precise population figure, which is one more reason we lean on the
shift in risk structure rather than on the headline percentage alone.

\paragraph{So what this implies for policy.}
Reading the collapse as broad-based and supply-side changes the policy prescription. If
the problem were concentrated household poverty, socioeconomically targeted transfers
would reach most at-risk students; because it is a flat, system-wide degradation of
schooling capacity, such transfers alone would miss them. The levers the evidence points
to instead are those that rebuild the delivery of instruction: retaining and paying
teachers to slow the brain drain, completing physical reconstruction of damaged schools,
funding sustained catch-up teaching, and addressing the two individual factors the model
ranks highest, math anxiety and the quality of teacher support. These are the same
system-wide levers that the shared, broad-based driver structure in
Section~\ref{sec:results} already implied, now grounded in the country's economic
history.
